%% ============================================================
%% Teil IV: Das Projekt – vollständige C++-Implementierung
%% Band 3 – Kapitel 12 bis 13
%% ============================================================

\part{Das Projekt}

% ─────────────────────────────────────────────────────────────
\chapter{Wir bauen motor\_node in C++ –
  Lifecycle, Serial und Service komplett}
% ─────────────────────────────────────────────────────────────

\begin{lernziele}
Du implementierst den vollständigen \texttt{motor\_node.cpp}
als rclcpp-Lifecycle-Node mit Serial-Kommunikation,
\texttt{/cmd\_vel}-Subscriber, \texttt{/motor/stopp}-Service
und korrekter State~Machine.
\end{lernziele}

\section{Situation}

In den vorangehenden Kapiteln wurden alle Teile einzeln entwickelt.
Jetzt kommen sie zusammen:
Lifecycle~Node + Serial-Port + Subscriber + Service = der vollständige \texttt{motor\_node}.

\begin{aufgabeBox}
Implementiere \texttt{motor\_node.cpp} vollständig.
Er soll nach \texttt{ros2 lifecycle set /motor activate}
auf \texttt{/cmd\_vel} reagieren und den Arduino ansteuern.
\end{aufgabeBox}

\section{motor\_node.hpp}

\begin{lstlisting}[style=cpp]
// include/robotik_uno_q_cpp/motor_node.hpp
#pragma once

#include <rclcpp/rclcpp.hpp>
#include <rclcpp_lifecycle/lifecycle_node.hpp>
#include <geometry_msgs/msg/twist.hpp>
#include <std_srvs/srv/trigger.hpp>

#include <atomic>
#include <memory>
#include <mutex>
#include <string>
#include <thread>

namespace robotik_uno_q {

class MotorNode : public rclcpp_lifecycle::LifecycleNode {
public:
    explicit MotorNode(const rclcpp::NodeOptions& opts = rclcpp::NodeOptions{});
    ~MotorNode() override;

    using CallbackReturn =
        rclcpp_lifecycle::node_interfaces::LifecycleNodeInterface::CallbackReturn;

    CallbackReturn on_configure(const rclcpp_lifecycle::State&) override;
    CallbackReturn on_activate(const rclcpp_lifecycle::State&) override;
    CallbackReturn on_deactivate(const rclcpp_lifecycle::State&) override;
    CallbackReturn on_cleanup(const rclcpp_lifecycle::State&) override;
    CallbackReturn on_shutdown(const rclcpp_lifecycle::State&) override;

private:
    // ─── Parameter ───────────────────────────────────────────────────────
    std::string port_{"/dev/arduino"};
    int         baudrate_{9600};
    float       vmax_{0.5f};   // m/s: maximale Translationsgeschwindigkeit
    float       spurbreite_{0.165f};  // Meter

    // ─── Serial ──────────────────────────────────────────────────────────
    int         fd_{-1};      // File Descriptor des Serial-Ports (POSIX)
    std::mutex  serial_mutex_;

    // ─── ROS2-Handles ────────────────────────────────────────────────────
    rclcpp::Subscription<geometry_msgs::msg::Twist>::SharedPtr sub_cmd_;
    rclcpp::Service<std_srvs::srv::Trigger>::SharedPtr         srv_stopp_;

    // ─── State ───────────────────────────────────────────────────────────
    std::atomic<bool> motoren_aktiv_{false};

    // ─── Private Methoden ────────────────────────────────────────────────
    bool _oeffne_serial(const std::string& port, int baud);
    void _schliesse_serial();
    bool _sende(const std::string& befehl);

    void _cb_cmd_vel(const geometry_msgs::msg::Twist::SharedPtr& msg);
    void _cb_stopp(
        const std_srvs::srv::Trigger::Request::SharedPtr&,
        std_srvs::srv::Trigger::Response::SharedPtr&);

    static std::pair<int, int> _twist_zu_pwm(
        float vx, float wz, float vmax, float spurbreite);
};

}  // namespace robotik_uno_q
\end{lstlisting}

\section{motor\_node.cpp -- vollständige Implementierung}

\begin{lstlisting}[style=cpp]
// src/motor_node.cpp
#include "robotik_uno_q_cpp/motor_node.hpp"

#include <fcntl.h>
#include <termios.h>
#include <unistd.h>

#include <algorithm>
#include <cmath>
#include <format>
#include <stdexcept>
#include <string>

namespace robotik_uno_q {

// ─── Konstruktor / Destruktor ─────────────────────────────────────────────

MotorNode::MotorNode(const rclcpp::NodeOptions& opts)
    : LifecycleNode("motor", opts)
{
    // Parameter deklarieren (Defaults wie in Python-Version)
    this->declare_parameter("port",      port_);
    this->declare_parameter("baudrate",  baudrate_);
    this->declare_parameter("vmax",      vmax_);
    this->declare_parameter("spurbreite", spurbreite_);

    // Service: immer erreichbar, unabhängig vom Lifecycle-State
    srv_stopp_ = this->create_service<std_srvs::srv::Trigger>(
        "/motor/stopp",
        [this](const std_srvs::srv::Trigger::Request::SharedPtr& req,
               std_srvs::srv::Trigger::Response::SharedPtr&      res) {
            this->_cb_stopp(req, res);
        });

    RCLCPP_INFO(get_logger(),
        "MotorNode erstellt. Service /motor/stopp aktiv. "
        "Warte auf 'configure'...");
}

MotorNode::~MotorNode() {
    _schliesse_serial();
}

// ─── Lifecycle-Callbacks ──────────────────────────────────────────────────

MotorNode::CallbackReturn
MotorNode::on_configure(const rclcpp_lifecycle::State&)
{
    // Parameter lesen
    port_       = this->get_parameter("port").as_string();
    baudrate_   = this->get_parameter("baudrate").as_int();
    vmax_       = static_cast<float>(this->get_parameter("vmax").as_double());
    spurbreite_ = static_cast<float>(this->get_parameter("spurbreite").as_double());

    RCLCPP_INFO(get_logger(),
        "Konfiguriere: port=%s, baudrate=%d", port_.c_str(), baudrate_);

    // Serial-Port öffnen
    if (!_oeffne_serial(port_, baudrate_)) {
        RCLCPP_ERROR(get_logger(), "Serial-Port nicht erreichbar: %s", port_.c_str());
        return CallbackReturn::FAILURE;
    }

    // Subscriber erstellen
    sub_cmd_ = this->create_subscription<geometry_msgs::msg::Twist>(
        "/cmd_vel", rclcpp::QoS(10).reliable(),
        [this](const geometry_msgs::msg::Twist::SharedPtr msg) {
            this->_cb_cmd_vel(msg);
        });

    RCLCPP_INFO(get_logger(), "Konfiguriert. Bereit fuer 'activate'.");
    return CallbackReturn::SUCCESS;
}

MotorNode::CallbackReturn
MotorNode::on_activate(const rclcpp_lifecycle::State&)
{
    // Stopp-Befehl senden, dann Motoren freigeben
    _sende("S\n");
    motoren_aktiv_.store(true);
    RCLCPP_INFO(get_logger(), "AKTIVIERT: Motoren bereit.");
    return CallbackReturn::SUCCESS;
}

MotorNode::CallbackReturn
MotorNode::on_deactivate(const rclcpp_lifecycle::State&)
{
    motoren_aktiv_.store(false);
    _sende("S\n");
    RCLCPP_INFO(get_logger(), "Deaktiviert: Motoren gestoppt.");
    return CallbackReturn::SUCCESS;
}

MotorNode::CallbackReturn
MotorNode::on_cleanup(const rclcpp_lifecycle::State&)
{
    sub_cmd_.reset();
    _schliesse_serial();
    RCLCPP_INFO(get_logger(), "Bereinigt.");
    return CallbackReturn::SUCCESS;
}

MotorNode::CallbackReturn
MotorNode::on_shutdown(const rclcpp_lifecycle::State&)
{
    motoren_aktiv_.store(false);
    _sende("S\n");
    _schliesse_serial();
    return CallbackReturn::SUCCESS;
}

// ─── Private Methoden ─────────────────────────────────────────────────────

bool MotorNode::_oeffne_serial(const std::string& port, int baud) {
    fd_ = ::open(port.c_str(), O_RDWR | O_NOCTTY | O_NDELAY);
    if (fd_ < 0) return false;

    termios tty{};
    ::tcgetattr(fd_, &tty);
    speed_t speed = (baud == 115200) ? B115200 : B9600;
    ::cfsetospeed(&tty, speed);
    ::cfsetispeed(&tty, speed);
    tty.c_cflag = CS8 | CREAD | CLOCAL;
    tty.c_iflag = IGNPAR;
    tty.c_oflag = 0;
    tty.c_lflag = 0;
    tty.c_cc[VMIN]  = 0;
    tty.c_cc[VTIME] = 10;  // 1s Timeout
    ::tcflush(fd_, TCIFLUSH);
    ::tcsetattr(fd_, TCSANOW, &tty);
    return true;
}

void MotorNode::_schliesse_serial() {
    std::lock_guard<std::mutex> lock(serial_mutex_);
    if (fd_ >= 0) {
        ::close(fd_);
        fd_ = -1;
    }
}

bool MotorNode::_sende(const std::string& befehl) {
    std::lock_guard<std::mutex> lock(serial_mutex_);
    if (fd_ < 0) return false;
    ssize_t n = ::write(fd_, befehl.c_str(), befehl.size());
    return n == static_cast<ssize_t>(befehl.size());
}

void MotorNode::_cb_cmd_vel(const geometry_msgs::msg::Twist::SharedPtr& msg) {
    if (!motoren_aktiv_.load()) return;

    auto [pwm_r, pwm_l] = _twist_zu_pwm(
        static_cast<float>(msg->linear.x),
        static_cast<float>(msg->angular.z),
        vmax_, spurbreite_);

    // Format aus Band 1: "V 150 150\n"
    std::string befehl = std::format("V {} {}\n", pwm_r, pwm_l);
    if (!_sende(befehl)) {
        RCLCPP_WARN(get_logger(), "Serial-Sendefehler: %s", befehl.c_str());
    }
}

void MotorNode::_cb_stopp(
    const std_srvs::srv::Trigger::Request::SharedPtr&,
    std_srvs::srv::Trigger::Response::SharedPtr& res)
{
    RCLCPP_WARN(get_logger(), "NOTBREMSE!");
    motoren_aktiv_.store(false);
    bool ok = _sende("S\n");
    res->success = ok;
    res->message = ok ? "Gestoppt." : "Serial-Fehler.";
}

std::pair<int, int> MotorNode::_twist_zu_pwm(
    float vx, float wz, float vmax, float spurbreite)
{
    float v_r = vx + (spurbreite / 2.0f) * wz;
    float v_l = vx - (spurbreite / 2.0f) * wz;
    auto  clamp_pwm = [vmax](float v) -> int {
        return static_cast<int>(std::clamp(v / vmax, -1.0f, 1.0f) * 255.0f);
    };
    return {clamp_pwm(v_r), clamp_pwm(v_l)};
}

}  // namespace robotik_uno_q

// ─── main ─────────────────────────────────────────────────────────────────

int main(int argc, char* argv[]) {
    rclcpp::init(argc, argv);
    auto node = std::make_shared<robotik_uno_q::MotorNode>();

    rclcpp::executors::StaticSingleThreadedExecutor executor;
    executor.add_node(node->get_node_base_interface());
    executor.spin();

    rclcpp::shutdown();
    return 0;
}
\end{lstlisting}

\begin{loesungBox}
\begin{lstlisting}[language=bash]
# Bauen
cd ~/ros2_ws
colcon build --packages-select robotik_uno_q_cpp
source install/setup.bash

# Starten
ros2 run robotik_uno_q_cpp motor_node

# Lifecycle-Transitionen
ros2 lifecycle set /motor configure
ros2 lifecycle set /motor activate

# Python-Node published, C++-Node empfängt
ros2 run robotik_uno_q navigation_node
\end{lstlisting}
\end{loesungBox}

\begin{projektBox}[robotik-uno-q -- Projektstand]
\textbf{Heute:} \texttt{motor\_node.cpp} vollständig implementiert.\\
\textbf{Nächstes Kapitel:} Hybrides Deployment -- Python und C++ zusammen.
\end{projektBox}

\begin{merksatzBox}
\textbf{Merksatz:}
POSIX-Serial (\texttt{open}/\texttt{write}/\texttt{close}) statt libserial
für maximale Kontrolle.
\texttt{std::format} (C++20) statt \texttt{sprintf} für String-Formatierung.
\texttt{std::pair<int,int>} für Rückgabe zweier Werte ohne Overhead.
Namespace \texttt{robotik\_uno\_q} für alle eigenen Klassen.
\end{merksatzBox}


% ─────────────────────────────────────────────────────────────
\chapter{Wir wollen Python und C++ zusammen –
  hybrides Deployment}
% ─────────────────────────────────────────────────────────────

\begin{lernziele}
Du startest den vollständigen Stack mit Python-Nodes (Band~2)
und C++-Nodes (Band~3) gemeinsam, aktualisierst die Launch-Datei
und führst einen vollständigen Systemtest durch.
\end{lernziele}

\section{Situation}

Band~2 hat Python-Nodes.
Band~3 hat C++-Nodes.
Jetzt laufen beide zusammen -- in einem Workspace,
einer Launch-Datei, einem ROS2-Domain.

\begin{aufgabeBox}
Aktualisiere \texttt{roboter.launch.py} aus Band~2 so,
dass \texttt{motor\_node} die C++-Version aus \texttt{robotik\_uno\_q\_cpp} startet.
Führe einen vollständigen Systemtest durch.
\end{aufgabeBox}

\section{Aktualisierte Launch-Datei}

\begin{lstlisting}
# stufe2_ros2/launch/roboter.launch.py -- Version 2.0
"""Gemischter Stack: Python arduino_bridge + navigation, C++ motor."""

from launch import LaunchDescription
from launch_ros.actions import Node, LifecycleNode

PYTHON_PKG = 'robotik_uno_q'
CPP_PKG    = 'robotik_uno_q_cpp'
CONFIG     = '...'  # wie in Band 2

def generate_launch_description() -> LaunchDescription:
    arduino_bridge = Node(
        package=PYTHON_PKG,
        executable='arduino_bridge_node',
        name='arduino_bridge',
        parameters=[CONFIG],
        output='screen',
    )

    navigation = Node(
        package=PYTHON_PKG,
        executable='navigation_node',
        name='navigation',
        parameters=[CONFIG],
        output='screen',
    )

    # C++ motor_node als LifecycleNode starten
    motor = LifecycleNode(
        package=CPP_PKG,
        executable='motor_node',
        name='motor',
        namespace='',
        parameters=[CONFIG],
        output='screen',
    )

    return LaunchDescription([arduino_bridge, navigation, motor])
\end{lstlisting}

\section{Lifecycle-Transitionen in der Launch-Datei}

\begin{lstlisting}
from launch_ros.actions import LifecycleNode
from launch.actions import EmitEvent, RegisterEventHandler
from launch_ros.events.lifecycle import ChangeState
from launch_ros.event_handlers import OnStateTransition
import lifecycle_msgs.msg

def generate_launch_description():
    motor = LifecycleNode(...)

    # Nach configure automatisch activate
    aktiviere = RegisterEventHandler(
        OnStateTransition(
            target_lifecycle_node=motor,
            goal_state='inactive',
            entities=[
                EmitEvent(event=ChangeState(
                    lifecycle_node_matcher=...,
                    transition_id=lifecycle_msgs.msg.Transition.TRANSITION_ACTIVATE
                ))
            ]
        )
    )
    return LaunchDescription([motor, aktiviere])
\end{lstlisting}

\section{Vollständiger Systemtest}

\begin{loesungBox}
\begin{lstlisting}[language=bash]
# Beide Pakete bauen
cd ~/ros2_ws
colcon build
source install/setup.bash

# Stack starten (Arduino angesteckt!)
ros2 launch robotik_uno_q roboter.launch.py

# In separaten Terminals prüfen:

# 1. Alle Nodes aktiv?
ros2 node list
# /arduino_bridge  /navigation  /motor

# 2. Topics fließen?
ros2 topic hz /sensor/rechts     # ~10 Hz
ros2 topic hz /cmd_vel           # ~10 Hz

# 3. Lifecycle-State des C++-Nodes
ros2 lifecycle get /motor
# - State: active [3]

# 4. Notbremse testen
ros2 service call /motor/stopp std_srvs/srv/Trigger
# response: success=True, message='Gestoppt.'

# 5. Sprachen des Node-Graphs überprüfen
ros2 node info /arduino_bridge   # Python
ros2 node info /motor            # C++

# 6. Roboter fährt autonom -- Hindernis vor den Sensor halten!
ros2 topic echo /cmd_vel --once
# linear: x: 0.0, angular: z: 0.0   <- Stopp bei <25 cm
\end{lstlisting}
\end{loesungBox}

\begin{projektBox}[robotik-uno-q -- Projektstand nach Band~3]
\begin{tabular}{lll}
\hline
\textbf{Node} & \textbf{Sprache} & \textbf{Aufgabe} \\
\hline
\texttt{arduino\_bridge} & Python~3 / rclpy & Serial lesen, Range publishen \\
\texttt{navigation}      & Python~3 / rclpy & cmd\_vel entscheiden \\
\texttt{motor}           & C++17 / rclcpp  & Lifecycle, Serial, Notbremse \\
\hline
\end{tabular}

\bigskip
\textbf{Nächster Band:} ROS2 vertieft -- SLAM, Nav2, micro-ROS.
Der Roboter bekommt eine Karte, plant Routen
und navigiert vollständig autonom.
\end{projektBox}

\begin{merksatzBox}
\textbf{Merksatz:}
Python und C++ coexistieren im gleichen ROS2-Stack.
Sprachunterschiede sind für andere Nodes unsichtbar --
sie sehen nur Topics, Services und Lifecycle-States.
Das ist die Stärke von ROS2: sauber getrennte Verantwortlichkeiten,
sprach- und prozessübergreifend.
\end{merksatzBox}

\begin{robotikBox}
Das Hybridmodell ist in der Industrie Standard:
Nav2 (C++), ros2\_control (C++),
Anwendungs-Nodes und Testskripte (Python).
Wer beide Sprachen beherrscht und weiß, wann welche passt,
kann jeden ROS2-Stack lesen, debuggen und erweitern.
Das ist das Ziel dieser zwei Bände.
\end{robotikBox}
